\section{Open Problems}
\label{sec:openproblems}

In this section, we highlight open problems arising from user-centric personalization, long-horizon interaction and credit assignment, world-model-based reasoning, multi-agent collaboration and training, latent internal reasoning, and the governance of agentic systems operating autonomously in real-world environments.


\subsection{User-centric Agentic Reasoning and Personalization}
User-centric agentic reasoning \cite{qian2025userrl,li2025hello} refers to an agent’s ability to tailor its reasoning and actions to a specific individual user by modeling user characteristics, preferences, and interaction history over time. Rather than optimizing a fixed, task-defined objective, a user-centric agent treats the user as part of the environment and continuously adapts its strategy through extended, multi-turn interaction. This requires the agent to dynamically infer evolving user intent, accommodate changes in goals and behavior styles, and adjust decisions based on explicit or implicit user feedback as the dialogue progresses. Crucially, user-centric agentic reasoning involves balancing short-term task rewards with long-term user experience, satisfaction, and trust, which introduces non-stationary objectives and long-horizon credit assignment challenges beyond conventional agentic reasoning settings.

\subsection{Long-horizon Agentic Reasoning from Extended Interaction}
A central open challenge in agentic reasoning is robust long-horizon planning and credit assignment across extended interactions. While methods such as ReAct and Tree of Thought improve short-horizon reasoning \citep{yao2023react,yao2023tree}, errors still compound rapidly in long tasks, as illustrated by embodied agents like Voyager \citep{wang2023voyager}. RL-trained agents such as WebRL and Agent-R1 improve performance in realistic environments but rely on heavily engineered, domain-specific rewards and largely treat episodes independently \citep{qi2024webrl,wei2025webagent}. More recent process-aware approaches attempt to construct finer-grained credit signals \citep{xiang2024retrospex,yan2025memoryr1,ji2025tree}, yet remain environment-specific. A core open problem is how to assign credit across tokens, tool calls, skills, and memory updates, and to generalize such learning across a long sequence of episodes and tasks.

\subsection{Agentic Reasoning with World Models}
World-model-based agents \cite{carbonneaux2025cwm, zhang2025agent2} aim to mitigate myopic reasoning by enabling internal simulation and lookahead. Model-based RL systems such as DreamerV3 demonstrate the effectiveness of imagined rollouts for long-horizon control \citep{hafner2023mastering}, while recent LLM-based agents adapt world models to web, code, and GUI environments \citep{tang2024worldcoder,carbonneaux2025cwm,chae2024web,luo2025vimo}. However, current designs rely on ad hoc representations and are typically trained on short-horizon or environment-specific data, raising concerns about calibration and generalization. Only a few works explore co-evolving world models and agents over long time scales \citep{fang2025webevolver,deng2025simura}. An open problem is how to jointly train, update, and evaluate world models in non-stationary environments, and how to assess their causal impact on downstream planning reliability.

\subsection{Multi-agent Collaborative Reasoning and Training}
Multi-agent collaboration has emerged as a powerful paradigm for scaling agentic reasoning through role specialization and division of labor \citep{li2023camel,chen2023agentverse,wu2024autogen}. While debate- and role-based systems often outperform single agents, most collaboration structures are still manually designed. Recent multi-agent RL approaches begin to treat collaboration itself as a trainable skill \citep{liu2025llm,park2025maporl,ma2024coevolving}, but credit assignment at the group level remains poorly understood. Scaling to larger agent populations further introduces challenges in topology adaptation, coordination overhead, and safety \citep{qian2024scaling,florian2025agentsnet, zhang2024knowledge}. A key open problem is how to learn adaptive, interpretable collaboration policies that remain robust under partial observability and adversarial conditions.


\subsection{Latent Agentic Reasoning}
Latent agentic reasoning \cite{du2025enabling,fu2025cache,latentmas} explores performing planning, decision-making and collaboration in internal latent spaces rather than explicit natural language or symbolic traces. Recent work suggests that latent reasoning can improve efficiency and scalability, but at the cost of reduced interpretability and controllability. In agentic settings, this raises additional challenges, including how to align latent reasoning with external objectives, tools, agents and memory systems. Diagnosing failures becomes particularly difficult when intermediate reasoning steps are not externally observable. An open problem is how to design learning objectives, probing methods, and evaluation benchmarks that make latent agentic reasoning both effective and auditable.


\subsection{Governance of Agentic Reasoning}
Governance is a cross-cutting challenge for agentic reasoning systems that act autonomously over tools, environments, and other agents. Beyond standard LLM safety issues, agentic systems introduce new risks due to long-horizon planning, persistent memory, and real-world action execution \citep{wang2025comprehensive}. Failures may arise from interactions across time and components, making attribution and auditing difficult. Existing benchmarks and guardrails mainly focus on short-horizon behaviors \citep{zhen2024guardagent,tongxin2024rjudge}, leaving planning-time failures and multi-agent dynamics underexplored. A central open problem is to develop governance frameworks that jointly address model-level alignment, agent-level policies, and ecosystem-level interactions under realistic deployment conditions.